% =====================================================================
% What an SO(3) orientation flow can and cannot learn in
% molecular-crystal structure prediction
%
% Formatted for npj Computational Materials (Nature Portfolio).
% For submission, replace the preamble with the official Springer Nature
% template:  \documentclass[pdflatex,sn-nature]{sn-jnl}  and move the
% metadata into the SN \author/\affil macros. The section structure,
% numbered references, and back-matter below already follow npj policy.
% Compiles standalone with pdflatex + bibtex (natbib / unsrtnat).
% =====================================================================
\documentclass[11pt]{article}

\usepackage[margin=1in]{geometry}
\usepackage{amsmath,amssymb}
\usepackage{graphicx}
\usepackage{booktabs}
\usepackage{microtype}
\usepackage[numbers,super,sort&compress]{natbib}
\usepackage[colorlinks=true,allcolors=blue]{hyperref}
\usepackage{siunitx}
\usepackage{authblk}

\renewcommand{\abstractname}{\vspace{-\baselineskip}}
\setlength{\parskip}{0.2em}

\title{\bfseries What an SO(3) orientation flow can and cannot learn in
molecular-crystal structure prediction}

\author[1]{Frank Cai\thanks{\texttt{cai495@purdue.edu}}}
\affil[1]{Purdue University, West Lafayette, IN, USA}
\date{}

\begin{document}
\maketitle

\begin{abstract}
\noindent
Rigid-body factorization is an appealing route to molecular-crystal structure
prediction: freezing each molecule's conformer reduces a crystal to a lattice,
a set of fractional centroids on the torus $\mathbb{T}^3$, and a per-molecule
orientation on $\mathrm{SO}(3)$, a low-dimensional target amenable to fast
flow-matching samplers. Whether the orientation field is actually learnable on
real molecular crystals has not been examined. We introduce SymMC-Flow, a
space-group-conditioned rigid-body flow on
$\text{lattice}\times\mathbb{T}^3\times\mathrm{SO}(3)$, and use it to
characterize the orientation field on $1{,}127$ crystals exported from the
Cambridge Structural Database. The lattice and centroid flows learn normally,
but the absolute per-molecule orientation flow remains at its predict-zero
floor. We trace this to a decomposition of the orientation target,
$R_m = \mathrm{rot}(g_m)\cdot R_{\text{asym}}$, into a space-group-determined
relative rotation between symmetry copies and the asymmetric unit's free
orientation in the cell. The free part is gauge-arbitrary and unlearnable; the
relative part is learnable. Re-gauging to cancel the free part lifts the
held-out non-reference orientation loss by $27\%$ (versus $\sim\!0\%$ for the
absolute target) and, in an orientation-isolated evaluation, reconstructs
$16.8\%$ of held-out multi-copy packings exactly under StructureMatcher against
$0\%$ for the predict-floor baseline. Conditioning on a discrete space-group
coset identity raises the learnable signal to $51\%$, showing the residual
ceiling is set by the difficulty of \emph{inferring} the symmetry operation
from packing rather than by the network's capacity to represent the rotation.
The result delineates exactly what a rigid-body $\mathrm{SO}(3)$ flow can and
cannot contribute to molecular-crystal generation.
\end{abstract}

\noindent\textbf{Keywords:} crystal structure prediction; flow matching;
molecular crystals; rigid-body factorization; $\mathrm{SO}(3)$; space-group
symmetry

% =====================================================================
\section{Introduction}
% =====================================================================

Crystal structure prediction (CSP) asks for the stable periodic arrangement of
a given chemical system, a problem whose combinatorial landscape has
historically demanded large ensembles of density-functional relaxations.
Generative models have emerged as a fast alternative that learns the
distribution of stable structures directly. The progression from the crystal
diffusion variational autoencoder\cite{xie2022cdvae} to joint equivariant
diffusion over lattices and fractional coordinates\cite{jiao2023diffcsp} and to
Riemannian flow matching for materials\cite{miller2024flowmm} has steadily
lowered both sample cost and error on inorganic benchmarks, where each atom is
an independent site.

Molecular crystals are harder. A molecular crystal couples flexible
intramolecular degrees of freedom to the rigid-body packing of whole molecules,
and the most common organic space groups place several symmetry-related copies
of a molecule in each cell, a regime where packing search has long been a
bottleneck in community blind tests\cite{price2023blindtest7} and a target for
data-efficient learned models\cite{obrien2021dataefficient}. Two design choices
have emerged. All-atom
generators predict every atomic coordinate jointly: OXtal\cite{jin2025oxtal}
learns the conditional distribution over conformation and packing with a
lattice-free diffusion model, and PackFlow\cite{subramanian2026packflow} flows
heavy-atom Cartesian coordinates and lattice parameters with a reinforcement
learning alignment step. Rigid-body generators instead treat building blocks as
rigid and predict only their pose. MOFFlow\cite{kim2025mofflow} demonstrates
this for metal--organic frameworks, casting metal nodes and organic linkers as
rigid bodies and applying $\mathrm{SE}(3)$ flow matching to their
roto-translations, which mirrors generative modeling on the rotation
group\cite{leach2022so3} and the frame-based generation developed for protein
backbones\cite{yim2023se3diffusion,yim2023frameflow}.

Rigid-body factorization is attractive for molecular crystals: it collapses
dozens of atomic coordinates per molecule to a single pose, exploits the
modularity of the cell, and supports the short sampling trajectories
characteristic of flow matching\cite{lipman2023flowmatching,chen2024riemannian}.
It also introduces a per-molecule orientation variable on the rotation group
$\mathrm{SO}(3)$. Whether that orientation field is learnable on real molecular
crystals is an open question. MOFFlow reports a working method but does not
analyze the orientation field in isolation, and the all-atom approaches sidestep
orientation altogether. No prior work asks what a rigid-body $\mathrm{SO}(3)$
orientation flow actually learns on molecular crystals, or where it fails.

This paper answers that question. We build SymMC-Flow, a
space-group-conditioned rigid-body flow on
$\text{lattice}\times\mathbb{T}^3\times\mathrm{SO}(3)$, and train it on a
corpus of $1{,}127$ molecular crystals exported from the Cambridge Structural
Database (CSD). The lattice and centroid flows learn; the absolute per-molecule
orientation flow does not, sitting at the predict-zero floor across two corpus
sizes and after a gauge correction. We explain this with a decomposition of the
orientation target and a chain of controlled diagnostics, and we show that the
learnable component is real geometry, not a loss artifact.

Our contributions are:
\begin{itemize}
\itemsep0.1em
\item A rigid-body $\mathrm{SO}(3)$ flow for molecular crystals with a
molecule-intrinsic orientation gauge that removes the spurious cross-crystal
reference frame.
\item The decomposition $R_m = \mathrm{rot}(g_m)\cdot R_{\text{asym}}$ and a
falsification chain (predict-zero floor, gauge fix, scaling, clean-packing
control, relative re-gauge) that isolates the learnable symmetry-relative
orientation from the gauge-free asymmetric-unit pose.
\item An orientation-isolated StructureMatcher metric showing the learned
relative orientation reconstructs $16.8\%$ of held-out multi-copy packings
exactly, against $0\%$ for the floor.
\item Evidence from discrete coset-identity conditioning that the residual
ceiling is inference-limited, not representational.
\end{itemize}

% =====================================================================
\section{Results}
% =====================================================================

\subsection{Rigid-body factorization and the flow}

SymMC-Flow represents a molecular crystal as a lattice $L\in\mathbb{R}^{3\times3}$,
and, for each molecule $m$, a fractional centroid $c_m\in\mathbb{T}^3$ and an
orientation $R_m\in\mathrm{SO}(3)$ acting on a shared, centered body-frame
conformer $\text{local}_m$. Atom positions follow
\begin{equation}
\text{cart}_{m,a} = c_m L + R_m\,\text{local}_{m,a},
\label{eq:factor}
\end{equation}
so every copy of a chemical species shares one conformer and differs only by its
pose. A conditional flow-matching objective\cite{lipman2023flowmatching} is
trained jointly on the three manifolds: a Euclidean path for a
volume-normalized lattice parametrization, a torus geodesic for centroids, and
an $\mathrm{SO}(3)$ geodesic for orientations, with the vector field conditioned
on the space group. An $\mathrm{E}(n)$-equivariant
encoder\cite{satorras2021egnn} embeds the conformer and pair-bias attention
mixes molecules within the cell. Figure~\ref{fig:schematic} summarizes the
factorization and the orientation decomposition developed below. Architecture,
manifold operations, and data are detailed in Methods.

\subsection{Lattice and centroid flows learn on real molecular crystals}

We first establish that the pipeline learns the packing degrees of freedom.
Table~\ref{tab:perhead} reports the held-out per-head flow-matching loss before
and after training on the CSD corpus. The lattice loss falls by $94\%$ and the
centroid loss by $30\%$, both converging within the first epoch. The absolute
orientation loss moves by only $3.3\%$ and shows no downward trend over
training. Lattice and centroid behave as on inorganic benchmarks; orientation
is the outlier.

\begin{table}[t]
\centering
\caption{Held-out per-head flow-matching loss, untrained $\rightarrow$ trained,
on the CSD molecular-crystal corpus (noised conditioning). Lattice and centroid
learn; the absolute orientation target stays at its predict-zero floor.}
\label{tab:perhead}
\begin{tabular}{lccc}
\toprule
head & untrained & trained & change \\
\midrule
lattice            & 1.24 & 0.08 & $+94\%$ \\
centroid           & 0.36 & 0.25 & $+30\%$ \\
orientation ($R_m$)& 5.37 & 5.18 & $+3.3\%$ \\
\bottomrule
\end{tabular}
\end{table}

\subsection{The absolute orientation flow sits at a predict-zero floor}

The orientation loss equals the expected squared tangent velocity
$\mathbb{E}\lVert u_R\rVert^2\approx 5.24$ that a network outputting zero would
incur. A trained orientation head that never beats this value has learned
nothing. Three controls rule out the obvious explanations. First, the floor is
unchanged when the corpus is scaled from $250$ to $1{,}127$ crystals, so it is
not data sparsity. Second, a molecule-intrinsic orientation gauge (Methods),
which removes the arbitrary cross-crystal reference frame that would otherwise
make $R_m$ a rotation relative to some molecule in an unrelated structure, is
necessary for a well-posed target but does not lift the floor. Third, a
clean-packing control feeds the field the true (un-noised) lattice and
centroids while still noising the orientation and keeping the orientation
target, which is the best case for any two-stage scheme that generates packing
first and orientation second. Orientation still moves by only $3.3\%$
(Table~\ref{tab:twobytwo}, top row), so noised conditioning is not the cause and
a two-stage architecture cannot help. The obstruction is in the target itself.

\subsection{The orientation target decomposes into a learnable and a free part}

Consider a molecule $m$ generated from an asymmetric-unit molecule by a
space-group operation with rotation part $\mathrm{rot}(g_m)$. Its absolute pose
factors as
\begin{equation}
R_m = \mathrm{rot}(g_m)\cdot R_{\text{asym}},
\label{eq:decomp}
\end{equation}
where $R_{\text{asym}}$ is the orientation of the asymmetric unit in the cell.
The first factor is determined by the space group and is shared structure across
crystals; the second is a free, per-crystal choice with no relation to
composition or packing, and it dominates the target. Because $R_{\text{asym}}$
is gauge-arbitrary, no model can predict it, which is exactly the floor.

To isolate the learnable factor we re-gauge each crystal so that the first copy
of every species is a reference with $R=I$ and the remaining copies carry only
the relative rotation $R'_m = R_m R_0^{-1}$, which cancels $R_{\text{asym}}$
(Methods). The transformation leaves the reconstructed crystal unchanged. To
avoid crediting a trivial ``predict identity'' solution, we split the
orientation loss into reference copies (target $I$) and non-reference copies
(the genuinely symmetry-determined targets), and report the latter.
Table~\ref{tab:twobytwo} gives the full $2\times2$ over target and conditioning.
The absolute target floors in both conditioning regimes. The relative target
descends by $27\%$ on held-out non-reference copies and generalizes, with an
identical result under noised and clean conditioning. The signal therefore
rides on the space group, which is conditioned directly, together with the
coarse centroid arrangement, and is robust to conditioning noise.

\begin{table}[t]
\centering
\caption{Held-out non-reference orientation loss, untrained $\rightarrow$
trained, across orientation target and conditioning. The absolute pose floors;
the relative (symmetry-determined) orientation is learnable and robust to
conditioning noise.}
\label{tab:twobytwo}
\begin{tabular}{lcc}
\toprule
target & noised conditioning & clean packing \\
\midrule
absolute $R_m$  & $5.37\!\rightarrow\!5.18$ ($+3.3\%$) & $5.35\!\rightarrow\!5.18$ ($+3.3\%$) \\
relative $R'_m$ & $5.37\!\rightarrow\!3.91$ ($+27.1\%$) & $5.37\!\rightarrow\!3.94$ ($+26.7\%$) \\
\bottomrule
\end{tabular}
\end{table}

\subsection{Relative orientation reconstructs real packings}

A drop in flow-matching loss need not correspond to correct geometry. We
therefore evaluate the learned orientation by reconstruction. Holding the true
lattice, centroid, and conformer fixed, we sample only the $\mathrm{SO}(3)$ field
and rebuild each crystal, then test whether it matches the ground truth under
pymatgen's StructureMatcher\cite{ong2013pymatgen}. The symmetry-induced relative
rotations are large, clustering near $180^\circ$ because the dominant organic
space groups relate copies by two-fold rotations and inversions, so a single
deterministic integration of the field collapses to the conditional mean and
underperforms. We therefore draw eight orientation samples per crystal and take
the best, the appropriate generative read for a multimodal target. Reporting a
controlled best-of-$k$ match alongside the geodesic orientation error, rather
than a single headline rate, follows recent cautions on match-rate
metrics\cite{martirossyan2025matching}.

Table~\ref{tab:matchrate} reports the outcome on $131$ held-out multi-copy
crystals. The oracle, which uses the true orientation, matches every structure,
confirming the metric's tolerances. The trained flow reconstructs $16.8\%$ of
held-out packings exactly, against $0\%$ for a random (predict-floor) field and
$1.5\%$ for the naive choice of leaving every copy un-rotated. Its mean
non-reference orientation error is $68^\circ$, below the $85^\circ$ of a random
field. The $27\%$ loss reduction is real reconstructed geometry.

\begin{table}[t]
\centering
\caption{Orientation-isolated reconstruction on $131$ held-out multi-copy
crystals (true lattice, centroid, and conformer; only $\mathrm{SO}(3)$ sampled;
best of eight draws). The learned field reconstructs packings the floor and
naive baselines cannot.}
\label{tab:matchrate}
\begin{tabular}{lcc}
\toprule
condition & match rate & non-ref.\ orient.\ error \\
\midrule
oracle (true orientation) & $100.0\%$ & $0^\circ$ \\
identity (naive, $R=I$)   & $1.5\%$   & $165^\circ$ \\
untrained (predict-floor) & $0.0\%$   & $85^\circ$ \\
trained                   & $16.8\%$  & $68^\circ$ \\
\bottomrule
\end{tabular}
\end{table}

\subsection{The ceiling is inference-limited, not representational}

Two experiments locate the remaining gap. Increasing model capacity and
training (a wider model trained longer) raises the non-reference signal from
$27\%$ to $33\%$, real but diminishing headroom that plateaus well above zero.
The decisive test conditions the field on a discrete space-group coset
identity. For a copy generated by a crystallographic operation, the relative
rotation $R'_m$ equals that operation's rotation part, so within a space group
the non-reference relative rotations fall into a small discrete set of cosets.
We cluster them into a per-space-group codebook and condition each molecule on
its coset index, a discrete label of the kind used in symmetry-conditioned
crystal generation\cite{jiao2024diffcsppp,wyckoffdiff2025,symmcd2025}, rather
than the continuous rotation itself. Table~\ref{tab:strengthen} shows that
coset conditioning raises the non-reference signal to $51\%$, nearly doubling
the learnable component relative to a paired control that reproduces the $27\%$
baseline on the identical corpus. Most of the ceiling is therefore the
difficulty of inferring which symmetry operation generated each copy from the
noised packing, not an inability of the $\mathrm{SO}(3)$ head to represent the
rotation. A residual remains, attributable to the free $R_{\text{asym}}$ and to
the axial gauge freedom of symmetric-top molecules.

\begin{table}[t]
\centering
\caption{Held-out non-reference orientation loss under model and conditioning
changes. Coset conditioning nearly doubles the learnable signal; the paired
control isolates the effect.}
\label{tab:strengthen}
\begin{tabular}{lcc}
\toprule
configuration & non-ref.\ loss & change \\
\midrule
baseline (relative target)        & $5.37\!\rightarrow\!3.91$ & $+27.1\%$ \\
\quad + capacity and training     & $5.35\!\rightarrow\!3.56$ & $+33.4\%$ \\
\quad + coset conditioning        & $5.36\!\rightarrow\!2.65$ & $+50.6\%$ \\
\quad coset control (no coset)    & $5.37\!\rightarrow\!3.91$ & $+27.1\%$ \\
\bottomrule
\end{tabular}
\end{table}

\subsection{Method validation on inorganic benchmarks}

For completeness we validated the SymMC-Flow architecture on standard
single-block inorganic benchmarks, where orientation is inactive. On carbon-24
and MP-20 the method is competitive with joint equivariant
diffusion\cite{jiao2023diffcsp} at roughly an order of magnitude fewer sampling
steps, winning some match cells and losing others. Full numbers and the
head-to-head protocol are given in Supplementary Information; these results
confirm the flow and sampler are sound and are not the focus of this study.

% =====================================================================
\section{Discussion}
% =====================================================================

The orientation field of a rigid-body molecular-crystal flow is neither fully
learnable nor uniformly broken. It splits into a space-group-determined relative
orientation that a flow learns and reconstructs, and a gauge-free
asymmetric-unit orientation that no model can predict. This split is a property
of the representation, not of the optimizer or the architecture, which is why it
survives gauge correction, data scaling, capacity increases, and clean
conditioning.

The decomposition reconciles the divergent design choices in the field. All-atom
generators\cite{jin2025oxtal,subramanian2026packflow} never expose
$R_{\text{asym}}$ as an explicit target and so do not encounter the floor,
trading the low dimensionality of the rigid-body representation for a target
without a gauge obstruction. MOFFlow\cite{kim2025mofflow} succeeds with a
rigid-body $\mathrm{SE}(3)$ flow in part because framework connectivity
constrains block orientations far more tightly than the free pose of a molecule
in a van der Waals packing, shrinking the unlearnable component. The present
analysis predicts that rigid-body $\mathrm{SO}(3)$ generation will be most
effective precisely where orientation is constrained by the surrounding network
or by high site symmetry, and least effective for low-symmetry organic packings.

On evaluation, the gap between the flat orientation loss and the $16.8\%$
reconstruction rate underscores why a single match-rate number is an unreliable
summary\cite{martirossyan2025matching}: the learnable signal is visible only
when the metric isolates orientation and samples the multimodal target. We
report the geodesic error and a controlled best-of-$k$ match for this reason.

This study has limits. A de-novo stable--unique--novel evaluation is not
meaningful here because the universal interatomic potentials used to judge
stability are trained on inorganic data\cite{deng2023chgnet} and do not transfer
to organic molecular crystals; we therefore evaluate by reconstruction against
known structures. Training is at CPU scale. The rigid-conformer gate excludes
flexible molecules by construction, and the conformer registry is shared across
crystals, so a species whose conformer differs between polymorphs beyond
tolerance is rejected. These choices keep the rigid-body benchmark honest but
narrow its coverage.

Two directions follow. A hybrid that predicts a rigid pose and then refines the
conformer would extend the approach to flexible molecules while keeping the
fast pose flow. Conditioning the orientation field on Wyckoff or coset
information at sampling time, suggested by the coset experiment, would convert
the inference-limited component into a usable signal whenever the symmetry
assignment is available, as it is in template-based generation.

% =====================================================================
\section{Methods}
% =====================================================================

\subsection{Rigid-body factorization and intrinsic gauge}
Each crystal is parsed into rigid molecular blocks by a covalent-bond graph with
periodic-boundary unwrapping. Molecules of one species are keyed by a
Weisfeiler--Lehman hash of the element-labelled bond graph and share a single
reference conformer; atom correspondence between a copy and its reference is the
minimum-RMSD graph isomorphism, so a molecule's own automorphisms inject no
spurious rotation. The reference conformer is rotated into a molecule-intrinsic
frame given by the principal axes of its gyration tensor, with per-axis sign
fixed by an element-weighted third moment and the frame forced right-handed.
This removes the arbitrary cross-crystal reference that otherwise makes $R_m$ a
rotation relative to an unrelated molecule, the gauge correction referenced in
Results. The pose $R_m$ of each copy is then the Kabsch rotation aligning it to
the intrinsic frame. Lattices are rebuilt from their six cell parameters to a
canonical frame.

\subsection{Manifolds and flow matching}
The lattice is flowed in a ten-dimensional parametrization of per-atom
log-volume and a determinant-one shape, with a straight-line conditional path.
Centroids follow torus geodesics; orientations follow $\mathrm{SO}(3)$
geodesics with the tangent velocity as target. We use conditional flow
matching\cite{lipman2023flowmatching} with optimal-transport coupling of the
exchangeable prior and data within each crystal\cite{pooladian2023multisample}.
Sampling integrates the field from the prior to the data manifold with a
deterministic Runge--Kutta scheme, using the $\mathrm{SO}(3)$ exponential map
for the orientation update; molecular-crystal samples used here take fifty
steps. The construction follows Riemannian flow
matching\cite{chen2024riemannian} and parallels frame-based generative models
for proteins\cite{yim2023se3diffusion,yim2023frameflow} and metal--organic
frameworks\cite{kim2025mofflow}.

\subsection{Network}
An $\mathrm{E}(n)$-equivariant graph network\cite{satorras2021egnn} encodes each
conformer into an invariant token. Tokens are augmented with the current
centroid, orientation, lattice, time, and a space-group embedding, then mixed by
pair-bias attention with pairwise features built from torus and Cartesian
displacements. Three heads predict the lattice, centroid, and orientation
velocities. The centroid field can be symmetrized over the space-group action.
For the coset experiment a per-molecule coset embedding is added to each token
when enabled.

\subsection{Relative gauge and coset codebook}
The relative gauge sets the first copy of each species to $R=I$ and replaces the
others by $R'_m = R_m R_0^{-1}$, with the shared conformer transformed by
$R_0$ so that reconstruction is unchanged; a reference mask records which copies
are trivial. The coset codebook clusters non-reference relative rotations within
each space group by geodesic distance into a shared set of discrete indices,
which become the per-molecule coset labels.

\subsection{Orientation-isolated match metric}
Reconstruction holds the true lattice, centroid, and conformer fixed and
integrates only the orientation field from a random $\mathrm{SO}(3)$ prior under
clean-packing conditioning. Reconstructed crystals are built with
Eq.~\eqref{eq:factor} and compared to the ground truth with StructureMatcher
(fractional and site tolerances as in pymatgen
defaults)\cite{ong2013pymatgen}. For each crystal we draw eight orientation
samples and count a match if any draw matches; we also report the
minimum mean geodesic error over draws on non-reference copies.

\subsection{Data}
The corpus is a seeded random sample of the CSD filtered to ordered, organic,
three-dimensional, non-polymeric molecular crystals within size and
$Z'$ bounds, yielding $1{,}127$ crystals and $5{,}048$ rigid blocks; $1{,}095$
contain at least two copies of a species and define the multi-copy subset used
for the relative and coset analyses. A fixed seed gives a $964/131$
train/validation split. CSD structures are not redistributable; a refcode
manifest and the export and analysis code reproduce the dataset from the CSD.

\section*{Data availability}
The CSD entries used cannot be redistributed under the CSD licence. A manifest
of refcodes and all derived, non-redistributable quantities required to
reproduce the analyses are available from the author on request, and the dataset
can be regenerated from the manifest with the released code against a licensed
CSD installation.

\section*{Code availability}
The SymMC-Flow implementation, the export and factorization pipeline, and the
diagnostic and evaluation scripts are available at
\url{https://github.com/fronkt/symmc-flow}.

\section*{Author contributions}
F.C. conceived the study, implemented the method and experiments, performed the
analysis, and wrote the manuscript.

\section*{Competing interests}
The author declares no competing interests.

\section*{Funding}
This work received no external funding or grant support. Any cloud compute used
for model training was funded personally by the author.

\section*{Use of AI tools}
The author used an AI coding and writing assistant to support software
development and manuscript drafting. All scientific claims, results, and the
final text were verified by the author, who takes full responsibility for the
content.

\bibliographystyle{unsrtnat}
\bibliography{references}

\end{document}
